\documentclass[11pt]{article}

\usepackage[margin=1in]{geometry}
\usepackage{amsmath,amssymb,amsthm,mathtools}
\usepackage{booktabs}
\usepackage{microtype}
\usepackage[hidelinks]{hyperref}

\newtheorem{proposition}{Proposition}
\newtheorem{remark}{Remark}
\newcommand{\E}{\mathbb E}
\newcommand{\Var}{\operatorname{Var}}
\newcommand{\Cov}{\operatorname{Cov}}
\newcommand{\tr}{\operatorname{tr}}
\newcommand{\diag}{\operatorname{diag}}
\newcommand{\R}{\mathbb R}

\title{Statistical Properties of Forecast-Selected Penalized Trend Estimators}
\author{Heriberto Espino Montelongo}
\date{\today}

\begin{document}

\maketitle

\begin{abstract}
Finite-difference penalized least-squares trend estimators are linear smoothers
when their smoothing parameter is fixed. In that regime, their expectation,
bias, covariance, effective degrees of freedom, influence matrix, and forecast
variance follow directly from the smoothing matrix. The statistical object
changes once smoothness is selected from the same time series by chronological
forecast validation. The resulting estimator,
\[
\widehat\tau_{\widehat S}
=
H_{\lambda(\widehat S)}y,
\]
is no longer a fixed linear smoother because the selected smoothness
\(\widehat S\) is itself random and data-dependent. This paper develops the
fixed-smoothness foundation needed to study that distinction and formulates the
main inferential questions created by forecast-driven smoothness selection:
selection-adjusted complexity, bias and variance, uncertainty for trend
functionals, uncertainty in forecast-optimal smoothness, and propagation of
selection uncertainty into trend forecasts. The purpose is not to rederive
classical linear-smoother theory as a novelty claim, but to identify precisely
what is known conditionally on a fixed smoother and what must still be
established for the forecast-selected estimator used in the companion
smoothness-selection paper.
\end{abstract}

\section{Motivation and separation from the companion papers}

The research program contains three distinct questions.

The companion smoothness-selection paper asks
\[
\widehat S
\in
\arg\min_{S\in[0,1]}F(S),
\]
where \(F(S)\) is a chronological rolling future-block forecast loss. Its
question is: \emph{which amount of smoothness should be selected for
forecasting?}

The numerical-methods paper assumes that objective is already defined and asks:
\emph{how can its relevant minima and global optimum be located reliably and
efficiently when the objective is multimodal?}

The present paper asks a third question:
\begin{quote}
Once smoothness has been selected from the data, what statistical properties
does the resulting trend estimator have, and how should its uncertainty be
quantified?
\end{quote}

For fixed \(S\), penalized least squares is a linear smoother. After
forecast-based selection,
\[
y
\longmapsto
\widehat S(y)
\longmapsto
H_{\lambda(\widehat S(y))}
\longmapsto
H_{\lambda(\widehat S(y))}y
\]
is a data-adaptive nonlinear mapping.

The central statistical boundary of this paper is therefore
\[
\boxed{
\text{fixed-smoother statistics}
\quad\text{versus}\quad
\text{statistics after forecast-driven smoothness selection}.
}
\]

\section{Model and notation}

Let
\[
y=(y_1,\ldots,y_N)^\top\in\R^N
\]
be an observed time series. For difference order \(d\), let
\[
D_d\in\R^{(N-d)\times N},
\qquad
Q=D_d^\top D_d.
\]
The baseline penalized least-squares estimator is
\[
\widehat\tau_\lambda
=
\arg\min_{\tau\in\R^N}
\left\{
\|y-\tau\|_2^2
+
\lambda\|D_d\tau\|_2^2
\right\},
\]
with solution
\[
\boxed{
\widehat\tau_\lambda
=
H_\lambda y,
\qquad
H_\lambda=(I+\lambda Q)^{-1}.
}
\]

The baseline model is deliberately the zero-drift, identity-weight case used by
the companion methodological papers. The broader Guerrero formulation with
nonzero difference drift or a non-identity observation-weight/covariance
structure is treated as an extension rather than silently mixed into the
baseline.

Because
\[
x^\top Qx
=
x^\top D_d^\top D_dx
=
\|D_dx\|_2^2
\ge0,
\]
\(Q\) is symmetric positive semidefinite. Under the standard finite-difference
operator, its nullity is \(d\). Thus, if
\[
Q
=
U\diag(
\underbrace{0,\ldots,0}_{d},
\delta_1,\ldots,\delta_{N-d}
)U^\top,
\qquad
\delta_j>0,
\]
then
\[
H_\lambda
=
U\diag\left(
\underbrace{1,\ldots,1}_{d},
\frac1{1+\lambda\delta_1},
\ldots,
\frac1{1+\lambda\delta_{N-d}}
\right)U^\top.
\]

The normalized smoothness coordinate is
\[
\boxed{
S(\lambda)
=
1-
\frac1{N-d}
\sum_{j=1}^{N-d}
\frac1{1+\lambda\delta_j}.
}
\]
It maps \(\lambda=0\) to \(S=0\), and the limiting
\(\lambda\to\infty\) model to \(S=1\).

\section{What is already known for fixed smoothness}

This section collects the statistical foundation that follows once
\(\lambda\), equivalently \(S\), is treated as fixed. These identities are
important for the present paper, but they should not be presented as new
without a literature-specific audit because closely related results are
standard in Whittaker--Henderson smoothing, smoothing splines, generalized
ridge/Tikhonov regularization, Hodrick--Prescott filtering, Gaussian
random-walk/GMRF models, and related state-space formulations.

\subsection{Linear-smoother structure}

For fixed \(\lambda\),
\[
\widehat\tau_\lambda=H_\lambda y
\]
is linear in \(y\). The matrix \(H_\lambda\) is symmetric positive
semidefinite, and every eigenvalue lies in \((0,1]\).

The penalized spectral directions are attenuated by
\[
\alpha_j(\lambda)=\frac1{1+\lambda\delta_j}.
\]
The null-space directions are not attenuated.

\begin{proposition}[Null-space reproduction]
If \(x\in\ker(D_d)\), then
\[
H_\lambda x=x
\]
for every \(\lambda\ge0\).
\end{proposition}

For the usual difference operator, \(\ker(D_d)\) contains discrete polynomial
sequences of degree at most \(d-1\).

\subsection{Effective degrees of freedom and normalized smoothness}

For a fixed linear smoother, define
\[
\operatorname{edf}(\lambda)=\tr(H_\lambda).
\]
Using the spectral representation,
\[
\tr(H_\lambda)
=
d
+
\sum_{j=1}^{N-d}
\frac1{1+\lambda\delta_j}.
\]
Therefore
\[
\boxed{
\operatorname{edf}(\lambda)
=
N-(N-d)S(\lambda).
}
\]

Equivalently,
\[
\boxed{
S
=
\frac{N-\operatorname{edf}}{N-d}.
}
\]

This identity gives a direct statistical interpretation of normalized
smoothness: \(S\) is the normalized fraction of penalizable effective
flexibility removed by the smoother. At the endpoints,
\[
S=0\Longleftrightarrow\operatorname{edf}=N,
\qquad
S=1\Longleftrightarrow\operatorname{edf}=d.
\]

A central later question is whether \(\tr(H_{\widehat\lambda})\) remains an
adequate complexity measure after \(\widehat\lambda\) is selected from the data.

\subsection{Expectation, bias, variance, covariance, and risk}

Assume first
\[
y=\tau+\varepsilon,
\qquad
\E(\varepsilon)=0,
\qquad
\Var(\varepsilon)=\sigma^2I.
\]
For fixed \(\lambda\),
\[
\E(\widehat\tau_\lambda)
=
H_\lambda\tau,
\]
and
\[
\boxed{
\operatorname{Bias}(\widehat\tau_\lambda)
=
(H_\lambda-I)\tau.
}
\]

Likewise,
\[
\boxed{
\Var(\widehat\tau_\lambda)
=
\sigma^2H_\lambda^2.
}
\]

More generally, if
\[
\Var(\varepsilon)=\Sigma,
\]
then
\[
\boxed{
\Var(\widehat\tau_\lambda)
=
H_\lambda\Sigma H_\lambda^\top.
}
\]

For squared Euclidean risk,
\[
\E\|\widehat\tau_\lambda-\tau\|_2^2
=
\|(H_\lambda-I)\tau\|_2^2
+
\sigma^2\tr(H_\lambda^2).
\]

\subsection{Linear functionals of the trend}

For any \(a\in\R^N\),
\[
a^\top\widehat\tau_\lambda
=
a^\top H_\lambda y,
\]
and
\[
\Var(a^\top\widehat\tau_\lambda)
=
a^\top H_\lambda\Sigma H_\lambda^\top a.
\]

This includes level, local slope, and local curvature. The difficult question
is whether intervals constructed conditionally on a selected \(\widehat S\)
have correct post-selection coverage.

\subsection{Influence and residuals}

For fixed \(\lambda\),
\[
\frac{\partial\widehat\tau_i}{\partial y_j}
=
(H_\lambda)_{ij}.
\]
Hence \(H_\lambda\) is an observation-to-fit influence map.

Residuals are
\[
e_\lambda=(I-H_\lambda)y,
\]
with
\[
\Var(e_\lambda)
=
(I-H_\lambda)\Sigma(I-H_\lambda)^\top.
\]

Unlike an OLS hat matrix, \(H_\lambda\) is generally not idempotent:
\[
H_\lambda^2\neq H_\lambda.
\]
Thus OLS residual degrees-of-freedom formulas should not be copied
mechanically.

\section{Forecasting with fixed smoothness}

Let \(G_{d,h}\) denote the fixed continuation matrix. Then
\[
\widehat z_\lambda
=
G_{d,h}H_\lambda y.
\]
Conditional on fixed \(\lambda\),
\[
\E(\widehat z_\lambda)
=
G_{d,h}H_\lambda\tau,
\]
and
\[
\boxed{
\Var(\widehat z_\lambda)
=
G_{d,h}H_\lambda\Sigma H_\lambda^\top G_{d,h}^\top.
}
\]

This is uncertainty in the estimated and extrapolated trend. A full predictive
distribution additionally requires future innovation or observation noise.

After smoothness selection there is a third source:
\[
\text{forecast uncertainty}
=
\text{trend-estimation uncertainty}
+
\text{future innovation uncertainty}
+
\text{smoothness-selection uncertainty},
\]
with the caveat that exact finite-sample variance need not decompose additively
without covariance terms.

\section{The new statistical object: forecast-selected smoothness}

The smoothness-CV paper selects
\[
\widehat S
\in
\arg\min_{S\in[0,1]}F(S).
\]
The resulting estimator is
\[
\boxed{
\widehat\tau_{\widehat S}
=
H_{\lambda(\widehat S)}y.
}
\]

The notation is deceptively similar to the fixed-\(S\) estimator, but
\[
\widehat S=\widehat S(y)
\]
is random. Consequently,
\[
y\mapsto H_{\lambda(\widehat S(y))}y
\]
is generally nonlinear.

This is the central source of nontriviality in the paper.

\subsection{Why fixed-smoother variance is not the whole answer}

Conditioning on \(\widehat S=s\), one can write
\[
H_{\lambda(s)}\Sigma H_{\lambda(s)}^\top.
\]
But if inference concerns the full procedure that first selects \(S\) and then
estimates the trend, the randomness of \(\widehat S\) must also be accounted
for.

The same warning applies to effective degrees of freedom, confidence bands,
tests of slope or curvature, and forecast uncertainty.

\section{What we want to know}

The following are research targets, not established novelty claims.

\subsection{Sampling behavior of selected smoothness}

We want to characterize
\[
\widehat S=\arg\min_S F(S)
\]
as a random quantity:
\begin{itemize}
\item its sampling or stability distribution;
\item dependence on horizon, sample length, residual dependence, noise scale,
and latent roughness;
\item uncertainty when \(F(S)\) is multimodal;
\item confidence sets for a clearly defined population/oracle
forecast-optimal \(S^\star\).
\end{itemize}

\subsection{Selection-adjusted effective degrees of freedom}

For fixed \(\lambda\),
\[
\operatorname{edf}=\tr(H_\lambda).
\]
For
\[
m(y)=H_{\lambda(\widehat S(y))}y,
\]
we want an effective complexity measure for the full adaptive mapping rather
than assuming
\[
\operatorname{edf}_{\rm adaptive}
=
\tr(H_{\widehat\lambda}).
\]

\subsection{Bias and variance after selection}

We want to understand
\[
\E\left[H_{\lambda(\widehat S)}y\right]
\]
and
\[
\Var\left(H_{\lambda(\widehat S)}y\right),
\]
not only their conditional fixed-\(S\) analogues.

Candidate approaches include exact results in simplified Gaussian settings,
local expansions around stable interior optima, bootstrap approximations, and
simulation-based coverage studies.

\subsection{Inference for level, slope, and curvature}

We want selection-aware inference for quantities such as
\[
e_t^\top H_{\lambda(\widehat S)}y,
\qquad
\Delta H_{\lambda(\widehat S)}y,
\qquad
\Delta^2H_{\lambda(\widehat S)}y.
\]

Scientifically useful targets include level, local slope, sign of slope,
curvature, turning-point location, and terminal slope/curvature used by
extrapolation.

\subsection{Propagation of smoothness uncertainty into forecasts}

For fixed \(\lambda\),
\[
\widehat z_\lambda=GH_\lambda y.
\]
Since
\[
\frac{\partial H_\lambda}{\partial\lambda}
=
-H_\lambda QH_\lambda,
\]
we have
\[
\frac{\partial\widehat z_\lambda}{\partial\lambda}
=
-GH_\lambda QH_\lambda y.
\]
Because \(S\) is monotone in \(\lambda\),
\[
\boxed{
\frac{\partial\widehat z}{\partial S}
=
-\frac{GH_\lambda QH_\lambda y}{S'(\lambda)}.
}
\]

This suggests a local delta-method term involving
\[
\left(
\frac{\partial\widehat z}{\partial S}
\right)
\Var(\widehat S)
\left(
\frac{\partial\widehat z}{\partial S}
\right)^\top,
\]
but its validity must be investigated. It may be poor near boundaries, flat
objectives, or switches between local minima.

\subsection{Multimodality as a statistical issue}

The numerical-methods paper treats multiple minima as an optimization problem.
Here the same phenomenon is statistical.

If two minima \(S_1\) and \(S_2\) have similar forecast loss but imply
substantially different effective degrees of freedom, terminal derivatives, or
forecast geometry, small perturbations may move the fitted procedure between
different smoothing regimes.

We want to characterize competing minima by
\[
\operatorname{edf},
\quad
\text{roughness},
\quad
\text{bias/variance},
\quad
\text{spectral attenuation},
\quad
\text{terminal derivatives},
\quad
\text{forecast geometry}.
\]

\section{A local expansion around a stable selected optimum}

Suppose the relevant optimum is unique, interior, and locally stable:
\[
\widehat S=S^\star+\Delta_S.
\]
Then formally,
\[
\widehat\tau_{\widehat S}
\approx
\widehat\tau_{S^\star}
+
\left.
\frac{\partial\widehat\tau_S}{\partial S}
\right|_{S^\star}
\Delta_S.
\]

Because
\[
\frac{\partial\widehat\tau}{\partial\lambda}
=
-H_\lambda QH_\lambda y,
\]
\[
\boxed{
\frac{\partial\widehat\tau}{\partial S}
=
-\frac{H_\lambda QH_\lambda y}{S'(\lambda)}.
}
\]

This suggests a first-order decomposition into fixed-\(S^\star\) variability
and smoothness-selection variability. However, \(\widehat S\) and \(y\) are
generated from overlapping data, so covariance terms matter. A rigorous
treatment must derive rather than assume the decomposition.

The expansion is not expected to describe boundary solutions or local-minimum
switching well.

\section{What would constitute a contribution}

The following are foundations, not sufficient contributions by themselves:
\begin{itemize}
\item \(\operatorname{edf}=\tr(H)\);
\item fixed-\(\lambda\) bias and variance;
\item leverage from \(H_{ii}\);
\item ordinary fixed-\(\lambda\) Gaussian intervals;
\item the Bayesian interpretation of a quadratic penalty;
\item relationships with HP, ridge/Tikhonov, or Whittaker smoothing;
\item polynomial continuation implied by the null space.
\end{itemize}

A publishable contribution should center on one or more nontrivial consequences
of forecast-driven selection:
\begin{enumerate}
\item selection-adjusted degrees of freedom;
\item uncertainty and coverage after forecast-CV smoothness selection;
\item sampling/stability theory for \(\widehat S\);
\item treatment of multimodal smoothness-selection uncertainty;
\item forecast-risk decomposition including selection uncertainty;
\item inference for terminal level/slope/curvature after selection;
\item horizon-dependent extrapolation-risk theory.
\end{enumerate}

The final paper should choose one or two primary theorem and simulation targets
rather than claiming all of these at once.

\section{Relationship to the smoothness-CV paper}

The distinction is
\[
\boxed{
\text{smoothness-CV: }
y\longmapsto \widehat S
}
\]
versus
\[
\boxed{
\text{statistical-properties: }
(y,\widehat S)
\longmapsto
\widehat\tau_{\widehat S}
\longmapsto
\text{bias, variance, complexity, and uncertainty}.
}
\]

The smoothness-CV paper may use simple fixed-smoother properties to interpret
\(S\), especially
\[
\tr(H)=N-(N-d)S.
\]
It should not absorb the full post-selection inference problem.

\section{Relationship to the numerical-methods paper}

The numerical paper answers how to compute candidate minima of \(F(S)\). The
present paper treats the selected minimum as a random output of a data-dependent
procedure.

Thus:
\[
\text{criterion}
\longrightarrow
\text{numerical solution}
\longrightarrow
\text{statistical consequences}.
\]

Multimodality is the main bridge: numerically it means several minima must be
found; statistically it means the selected estimator may jump between
substantially different smoothing regimes.

\section{Planned simulation program}

The first simulations should be designed for inference rather than forecasting
performance. The latent trend should be known so that bias, variance, coverage,
and selection behavior can be measured.

Controlled factors should include:
\begin{itemize}
\item sample length;
\item forecast horizon;
\item observation-noise scale;
\item latent-trend roughness;
\item residual persistence;
\item distance of the optimum from \(S=0\) or \(S=1\);
\item unimodal versus multimodal forecast-CV geometry.
\end{itemize}

At minimum, compare:
\begin{enumerate}
\item an oracle fixed-\(S\) estimator;
\item a fixed but misspecified \(S\);
\item forecast-CV-selected \(\widehat S\);
\item conditional intervals that ignore selection;
\item one or more selection-aware uncertainty procedures.
\end{enumerate}

Primary endpoints should include bias, variance, mean-squared error, coverage,
interval width, effective-complexity estimates, the distribution/stability of
\(\widehat S\), and forecast coverage where relevant.

\section{Literature boundary}

A dedicated literature audit is required before any novelty claim. The closest
bodies of theory include Whittaker--Henderson smoothing, smoothing splines and
P-splines, generalized ridge/Tikhonov regularization, Hodrick--Prescott
filtering, Gaussian random walks and GMRFs, state-space trend models,
generalized degrees of freedom, data-adaptive smoothing, selective inference,
and uncertainty after tuning-parameter selection.

The repository already contains especially relevant material on Guerrero's
finite-difference PLS/controlled-smoothness formulation and modern
Whittaker--Henderson treatments. These make clear that fixed-smoother effective
degrees of freedom, smoothing bias, Bayesian uncertainty, and extrapolation
uncertainty cannot by themselves be claimed as new here.

The novelty audit should therefore be organized around the narrower question:
\begin{quote}
What statistical theory already exists for a smoother whose tuning parameter is
selected specifically by chronological forecast loss, especially when the
validation surface is dependent, horizon-specific, and potentially multimodal?
\end{quote}

\section{Current status}

Established within the repository:
\begin{itemize}
\item the finite-difference PLS estimator \(H_\lambda y\);
\item normalized smoothness \(S(\lambda)\);
\item \(\operatorname{edf}=N-(N-d)S\) for fixed smoothness;
\item first and second derivatives of \(H_\lambda\);
\item first and second derivatives of forecast loss;
\item the forecast operator \(GH_\lambda\);
\item empirical evidence that forecast-CV surfaces can be multimodal;
\item a companion numerical method for recovering relevant minima.
\end{itemize}

Not yet established:
\begin{itemize}
\item the distribution and uncertainty of \(\widehat S\);
\item selection-adjusted effective degrees of freedom;
\item bias and variance of the full selected procedure;
\item valid post-selection intervals for level/slope/curvature;
\item a rigorous decomposition of forecast uncertainty including smoothness
selection;
\item treatment of multimodal smoothness-selection uncertainty.
\end{itemize}

These open items, rather than the fixed-smoother formulas alone, define the
scientific target.

\section{Conclusion}

Finite-difference penalized trend estimation has a simple statistical
description when smoothness is fixed. That simplicity provides an exact
baseline against which the forecast-selected estimator can be measured. The
main statistical challenge begins when smoothness is selected from
chronological forecast performance. Then the fitted trend is no longer a fixed
linear smoother, and its complexity, bias, variance, interval coverage, and
forecast uncertainty must reflect the randomness and possible multimodality of
the selection step.

The intended contribution is therefore not a new derivation of classical
smoothing identities. It is a statistical theory and inferential framework for
the forecast-selected penalized trend estimator produced by the companion
smoothness-CV procedure.

\end{document}
